\section{Commercial network tools: spatial impacts, exposure designs and switchbacks}
\label{exec20:network-scope}

A branch intervention can affect customers outside the treated branch, while a
borrower's funding choice can propagate through guarantees and counterparties.
Commercial software already offers several relevant routes, but they answer
different questions. A spatial regression estimates a response under a supplied
weight matrix and model; an exposure experiment compares outcomes under an
assignment design; a switchback changes a population's treatment over time.
The proposal must compare these routes at the level of their targets, information
and assumptions. A graph interface or a feature called ``impact'' is not a
common causal estimand.

The selected primary documentation supports concrete profiles for Stata,
SAS/ETS, ArcGIS Pro and Statsig. Their strengths include supplied-network
estimation, spatial data preparation, diagnostic output, experimental assignment
and analysis services that our proposed library has not implemented. The
comparison therefore asks which original engines or services merit a public
prototype, and which semantics the proposed common interface must preserve.
It does not claim existing product parity.
\citep{exec20_network_stata_spregress,exec20_network_sas_latest_spatial,
exec20_network_arcgis_tool,exec20_network_statsig_types}

\subsection{Define the intervention before comparing an impact}
\label{exec20:network-target}

Let $Y$ be an $n$-vector of branch outcomes, $a$ a vector of intervention levels,
and $W$ a supplied matrix whose entry $W_{ij}$ weights the influence of unit $j$
on unit $i$. This receiver-row convention differs from the source-row adjacency
used by some discovery packages. For a stipulated linear structural equilibrium,
consider
\begin{equation}
 Y=\rho WY+\beta a+\gamma Wa+U,
 \qquad
 S(W)=\frac{\partial E[Y\mid\operatorname{do}(a)]}{\partial a^{\mathsf T}}
      =(I-\rho W)^{-1}(\beta I+\gamma W).
 \label{exec20:network-response}
\end{equation}
The second expression requires an invariant structural mechanism, an exogenous
intervention in $a$, a fixed relevant $W$, and an invertible feedback operator.
The sufficient stability restriction $|\rho|r(W)<1$, where $r(W)$ is the
spectral radius, also gives the convergent propagation expansion
$S=\sum_{k\geq0}\rho^kW^k(\beta I+\gamma W)$. A fitted conditional spatial
mean can have the same algebra without supplying those intervention premises.
An endogenous treatment, omitted common cause or treatment-dependent network
requires additional identification information. The equation is a project
derivation for the stated model, not a theorem about every vendor estimator.
\citep{exec20_network_reference}

The average own derivative is $D=\operatorname{tr}(S)/n$, the average total
derivative is $T=\boldsymbol{1}^{\mathsf T}S\boldsymbol{1}/n$, and the
average spillover derivative is $T-D$. These average quantities also differ from
the response to treating one selected branch $j$, which is the column $S e_j$.
Stata's \texttt{estat impact} documents the averages of own, off-diagonal and
total derivatives of its reduced-form mean, and exposes their Jacobians and
covariance matrices. ArcGIS Pro 3.5 instead documents ``simple impacts,'' with
the direct impact equal to the coefficient. A common result named
\texttt{direct\_effect} must preserve which definition was used.
\citep{exec20_network_stata_impacts,exec20_network_arcgis_method}

For an exact banking reference, set $\gamma=0$, $\beta=10$ dollars per
intervention unit, $\rho=1/2$, and
\begin{equation}
 W=\begin{pmatrix}0&1\\1&0\end{pmatrix},\qquad
 S=\begin{pmatrix}40/3&20/3\\20/3&40/3\end{pmatrix},\qquad
 (D,T-D,T)=(40/3,20/3,20).
 \label{exec20:network-feedback}
\end{equation}
At an assumed cost of USD 12 to intervene at branch 1, the raw coefficient
gives a margin of USD $-2$, the own-branch response gives USD $4/3$, and the
two-branch portfolio response gives USD 8. These are three target calculations,
not contradictory estimates of one quantity. If the actual tie from branch 1
to branch 2 is missing from the supplied matrix, the modeled portfolio response
becomes 10 and its margin becomes USD $-2$. Neither more optimizer iterations
nor a generic robust covariance option restores the omitted path. The example
uses an authored structural law and assumed costs; it is not observed bank or
commercial-runtime evidence. \citep{exec20_network_reference}

\subsection{Weights, identities and uncertainty are substantive choices}
\label{exec20:network-weights}

Stata distinguishes scalar normalization from row normalization. Dividing $W$
by a common scalar and changing spatial coefficients correspondingly can
preserve the model. Normalizing each row separately generally changes its
meaning. The manual also refuses a sample smaller than the matrix's unit set
unless \texttt{force} is requested, warning that removing units changes spatial
connections and normalization. Its maximum-likelihood robust covariance is
described for nonnormal independent and identically distributed innovations;
the heteroskedastic GS2SLS option relaxes identical distribution while retaining
independence. None of these choices supplies arbitrary residual network
dependence. \citep{exec20_network_stata_spregress}

The exact three-branch star
\begin{equation}
 W_\star=\begin{pmatrix}0&1&1\\1&0&0\\1&0&0\end{pmatrix},\qquad
 W_{\rm row}=\begin{pmatrix}0&1/2&1/2\\1&0&0\\1&0&0\end{pmatrix}
 \label{exec20:network-normalization}
\end{equation}
illustrates the distinction. With $\rho=1/4$, $\beta=10$ and $\gamma=0$, a unit
intervention at branch 1 has responses $(80/7,20/7,20/7)$ under $W_\star$ and
$(32/3,8/3,8/3)$ under $W_{\rm row}$. In contrast, replacing $W_\star$ by
$W_\star/2$ and $\rho$ by $1/2$ preserves the original response exactly.
The two row conventions describe different propagation assumptions. An agent
must not select one merely to match a package default.
\citep{exec20_network_reference}

SAS/ETS's selected SPATIALREG manual, labeled 2023.10 and dated August 2024,
documents a broad cross-sectional model family including SAR, spatial Durbin,
error and moving-average specifications. Its default row standardization
replaces missing weights with zeros; \texttt{NONORMALIZE} instead refuses missing
off-diagonal entries and treats diagonal entries as zero. A missing relationship
is not scientifically equivalent to a known absent relationship. The
\texttt{SPATIALID} field aligns observations with the matrix, but numeric
noninteger IDs are truncated for matching. The proposed interface should use
validated immutable identifiers and an explicit missing-tie policy before
dispatch, preserving all exclusions and conversions.
\citep{exec20_network_sas_latest_spatial}

The selected SAS SAR likelihood uses independent normal innovations. It offers
Hessian, outer-product and QML covariance choices, with Hessian as default and
additional restrictions when certain likelihood approximations are selected.
Stata's default impact delta method conditions on regressors; its unconditional
option includes their sampling variation and has sample-selection restrictions.
These documented distinctions are useful requirements for an adapter, but
do not by themselves prove identification or useful finite-sample coverage.
\citep{exec20_network_sas_latest_spatial,exec20_network_stata_impacts}

Full parameter covariance must also survive extraction. In the two-branch
reference, $T=\beta/(1-\rho)$ has gradient $(2,40)$ in the order $(\beta,\rho)$.
For the illustrative covariance
\begin{equation}
 V=\begin{pmatrix}1&1/20\\1/20&1/100\end{pmatrix},\qquad
 \nabla T^{\mathsf T}V\nabla T=28,
 \label{exec20:network-variance}
\end{equation}
whereas dropping the cross-covariance gives 20. The covariance is assumed and
positive definite; the calculation is an exact delta-method reference, not an
estimated standard error or coverage finding. It explains why a CLI or MCP
response containing only coefficients and individual standard errors cannot
generally support the intended propagated-value calculation.
\citep{exec20_network_reference}

Stata's spatial-panel route further requires strongly balanced panels. Its
fixed-effects formulation eliminates panel effects and cannot fit covariates
constant within a panel; random-effects formulations impose their own
distribution and independence restrictions. The displayed innovations are
independent across panels and time. This is a concrete panel capability worth
comparing, but it does not automatically qualify evolving guarantee networks or
dependent servicing histories. \citep{exec20_network_stata_panel}

ArcGIS Pro adds geographical neighborhood construction, diagnostic maps,
lag/error/combined model choices and \texttt{arcpy.stats.SAR}. Its documented
defaults and inferred bandwidths need to be retained with the result. The
method page describes robust errors and an AK-test-triggered spatial HAC option,
using a separate nearest-neighbor triangular-kernel matrix. It also documents a
30-percent connectivity refusal. That is a rule of the selected tool, not a
universal boundary for causal validity. Basic, Standard and Advanced availability
is listed, but the listing does not establish redistribution, embedding or
deployment rights. The selected prose is inspected; image-only equations and
flowcharts still need a separate original-source review.
\citep{exec20_network_arcgis_method,exec20_network_arcgis_tool}

\subsection{Exposure support is a property of the experiment}
\label{exec20:network-exposure}

For a different route, let $Z_i$ be treatment assignment and let
$E_i(Z,G)$ summarize the assignments of unit $i$'s neighbors in a declared
network $G$. An exposure response $Y_i(z,e)$ assumes that this summary retains
the relevant treatment interference. Even if the graph is measured perfectly,
the assignment design may never generate some required $(z,e)$ combinations.
For two connected borrowers with exactly one treated, the only possible
exposures are $(1,0)$ and $(0,1)$, each with probability $1/2$ for a given
borrower. The cells $(0,0)$ and $(1,1)$ have zero probability. Direct effects at
zero neighbor treatment cannot be recovered by weighting observations into a
cell the experiment never visits. \citep{exec20_network_reference}

Independent half-probability treatment assignments instead give each exposure
cell probability $1/4$. With known positive probabilities
$\pi_i(z,e)=\Pr\{Z_i=z,E_i=e\}$, an authored finite-population mean reference is
\begin{equation}
 \widehat\mu(z,e)=\frac1n\sum_{i=1}^n
 \frac{\mathbf{1}\{Z_i=z,E_i=e\}Y_i}{\pi_i(z,e)}.
 \label{exec20:network-ht}
\end{equation}
Under consistency for the declared exposure mapping and the known assignment
law, linearity of expectation yields the mean potential outcome for that cell.
Enumerating all four assignments for the stipulated law $Y_i(z,e)=5+2z+3e$
gives expected means 5, 8, 7 and 10 for $(0,0),(0,1),(1,0),(1,1)$ respectively.
The direct contrast at $e=0$ is 2, and the spillover contrast for untreated units
is 3. This verifies a point-mean identity. It supplies neither a useful variance
estimate from two borrowers nor a general confidence interval; joint exposure
probabilities, dependence and the sampling/design law remain necessary.
\citep{exec20_network_reference}

Statsig's selected cluster-unit documentation gives a different, concrete
service: assign at organization level and analyze user revenue using a specified
one-to-many organization/user relationship. Its ratio treatment accounts for
numerator-denominator covariance with the delta method. That can address the
distinction between assignment and analysis units, but does not establish that
all interactions lie within organizations, that treatment cannot alter the
denominator, or that the desired exposure cells have positive probability.
Those questions must remain in the proposed experiment contract.
\citep{exec20_network_statsig_units}

\subsection{Switchbacks change interference and time together}
\label{exec20:network-switchback}

Statsig's legacy technical documentation describes switchbacks in which
populations or regions switch treatment by time window. It documents event
attribution using exposure timestamps, unit IDs and an attribution window;
bucket-level and variant-level metrics; separate resampling of treatment and
control buckets 10,000 times; and percentile confidence intervals. Configuration
includes schedules, window lengths and burn-in/burn-out exclusions, with
randomized bucketing described as an advanced feature. These are substantive
capabilities for a public experimental comparison, not merely a network label.
\citep{exec20_network_statsig_types,exec20_network_statsig_switchback}

Changing everyone in a region together can reduce one form of simultaneous
cross-treatment interference while introducing an important time comparison.
For an authored law $Y_t(z)=100t+2z$, observing control at $t=0$ and treatment at
$t=1$ gives an observed difference of 102 although the immediate causal effect
is 2. With carryover $Y_t=2Z_t+4Z_{t-1}$ under an alternating schedule, treated
windows follow control and yield 2, while control windows follow treatment and
yield 4. The observed difference is $-2$, the immediate effect is 2, and the
sustained treatment effect is 6. Resampling the same contaminated windows cannot
create missing assignment or washout information. These are explicit
counterexamples to an unrestricted interpretation, not claims of an observed
Statsig failure. \citep{exec20_network_reference}

The prototype should therefore declare the policy target, randomization of
start phases and schedules, carryover horizon, attribution overlap and effective
resampling unit. Discarding exposures during burn periods is not by itself a
demonstration that all associated outcomes are free of earlier treatment.
Adjacent buckets or shared regional shocks require a justified dependence
treatment. The saved September 2025 documentation also lists profile limits,
including unavailable time-series and dimension breakdowns and unavailable CUPED
and sequential testing for that legacy switchback profile. Those statements
must not be promoted to claims about current commercial availability.
\citep{exec20_network_statsig_switchback}

\subsection{A comparison that can guide the first prototype}
\label{exec20:network-comparison}

For an initial equilibrium-model prototype, Stata and SAS provide documented
original engines with different model, covariance and weight policies. ArcGIS
offers a geographical workflow and a Python surface, with an impact convention
that needs explicit reconciliation. Statsig supplies an experimental alternative
when population/time assignment is the intended intervention. A rational first
comparison chooses a declared target and then asks which of these profiles
actually supports it; it does not rank unlike outputs by one accuracy number.
Our present implemented contribution is the exact authored references. Model
fitting, licensed execution, native covariance behavior and service integration
remain proposed capabilities.

The public prototype can begin with the two- and three-branch laws above,
generated assignments, missing-tie perturbations and explicit carryover. It
should compare target preservation, refusals, graph and sample identities,
uncertainty definitions and economic decisions before scale. Public empirical
datasets can test import, geography and workflow behavior, but their observed
network alone supplies no intervention truth. Actual runtime, memory, analyst
effort, licensing, data movement and upgrade costs must be measured rather than
inferred from documentation. Commercial adapters should preserve the original
engine where justified; porting a proprietary algorithm is neither necessary
nor supported by the present evidence.

Coverage remains incomplete for the exact editions and interference targets of
Optimizely, EViews, IBM SPSS and the commercial probabilistic-network families.
The requested Optimizely route returned a general homepage, which supports no
technical capability claim. Supplied Bayesian networks and ordinary pooled
regression must not stand in for an uninspected exposure estimator. Their next
steps are concrete: identify the exact procedure and version, acquire the
original technical specification and rights terms, and compare its target with
the same public reference. None is classified as lacking a feature merely
because its original has not yet been reviewed. These unresolved profiles and
unknown prices limit the market comparison while allowing the supported
prototype designs to proceed.
